\documentclass[10pt,a4paper]{amsart}
\usepackage{amsmath,amssymb,amsthm,graphicx}
% Retained minimal preamble for the Stacks Project chapter sites-cohomology.tex (Cohomology on Sites).
% Verbatim excerpt of the chapter preamble: Xy-pic setup (native preamble lines 14-18) and the native
% theorem declarations with the mathematical macros (native preamble lines 154-217). Copied mechanically;
% nothing added, removed or redefined. The wrapper class and the amsmath/amssymb/amsthm/graphicx base
% packages are supplied by the assembler, so the class/externaldocument part of the native preamble is omitted.
% For commutative diagrams we use Xy-pic
\usepackage[all]{xy}

% We use 2cell for 2-commutative diagrams.
\xyoption{2cell}
% Theorem environments.
%
\theoremstyle{plain}
\newtheorem{theorem}[subsection]{Theorem}
\newtheorem{proposition}[subsection]{Proposition}
\newtheorem{lemma}[subsection]{Lemma}

\theoremstyle{definition}
\newtheorem{definition}[subsection]{Definition}
\newtheorem{example}[subsection]{Example}
\newtheorem{exercise}[subsection]{Exercise}
\newtheorem{situation}[subsection]{Situation}

\theoremstyle{remark}
\newtheorem{remark}[subsection]{Remark}
\newtheorem{remarks}[subsection]{Remarks}

\numberwithin{equation}{subsection}

% Macros
%
\def\lim{\mathop{\mathrm{lim}}\nolimits}
\def\colim{\mathop{\mathrm{colim}}\nolimits}
\def\Spec{\mathop{\mathrm{Spec}}}
\def\Hom{\mathop{\mathrm{Hom}}\nolimits}
\def\Ext{\mathop{\mathrm{Ext}}\nolimits}
\def\SheafHom{\mathop{\mathcal{H}\!\mathit{om}}\nolimits}
\def\SheafExt{\mathop{\mathcal{E}\!\mathit{xt}}\nolimits}
\def\Sch{\mathit{Sch}}
\def\Mor{\mathop{\mathrm{Mor}}\nolimits}
\def\Ob{\mathop{\mathrm{Ob}}\nolimits}
\def\Sh{\mathop{\mathit{Sh}}\nolimits}
\def\NL{\mathop{N\!L}\nolimits}
\def\CH{\mathop{\mathrm{CH}}\nolimits}
\def\proetale{{pro\text{-}\acute{e}tale}}
\def\etale{{\acute{e}tale}}
\def\QCoh{\mathit{QCoh}}
\def\Ker{\mathop{\mathrm{Ker}}}
\def\Im{\mathop{\mathrm{Im}}}
\def\Coker{\mathop{\mathrm{Coker}}}
\def\Coim{\mathop{\mathrm{Coim}}}

% Boxtimes
%
\DeclareMathSymbol{\boxtimes}{\mathbin}{AMSa}{"02}

%
% Macros for moduli stacks/spaces
%
\def\QCohstack{\mathcal{QC}\!\mathit{oh}}
\def\Cohstack{\mathcal{C}\!\mathit{oh}}
\def\Spacesstack{\mathcal{S}\!\mathit{paces}}
\def\Quotfunctor{\mathrm{Quot}}
\def\Hilbfunctor{\mathrm{Hilb}}
\def\Curvesstack{\mathcal{C}\!\mathit{urves}}
\def\Polarizedstack{\mathcal{P}\!\mathit{olarized}}
\def\Complexesstack{\mathcal{C}\!\mathit{omplexes}}
% \Pic is the operator that assigns to X its picard group, usage \Pic(X)
% \Picardstack_{X/B} denotes the Picard stack of X over B
% \Picardfunctor_{X/B} denotes the Picard functor of X over B
\def\Pic{\mathop{\mathrm{Pic}}\nolimits}
\def\Picardstack{\mathcal{P}\!\mathit{ic}}
\def\Picardfunctor{\mathrm{Pic}}
\def\Deformationcategory{\mathcal{D}\!\mathit{ef}}
\usepackage{cleveref}
\crefname{originalreference}{original reference}{original references}
\Crefname{originalreference}{Original reference}{Original references}
\makeatletter
\newlabel{lemma-local-actual}{{08C9}{source}{}{sourcepointer-1}{}}
\newlabel{lemma-local-actual@cref}{{[originalreference][0][]08C9}{[1][0][]source}{}{}{}}
\makeatother
\crefname{originalsection}{original section}{original sections}
\Crefname{originalsection}{Original section}{Original sections}
\makeatletter
\newlabel{section-internal-hom}{{017G}{source}{}{sourcepointer-2}{}}
\newlabel{section-internal-hom@cref}{{[originalsection][0][]017G}{[1][0][]source}{}{}{}}
\makeatother
\title{The internal hom of a strictly perfect complex with any complex is represented by the termwise internal hom complex}
\author{The Stacks Project Authors}
\date{stacks-project:sites-cohomology.tex:local-snapshot-20261011; GNU FDL 1.2 or later}
\begin{document}
\maketitle
\noindent\textit{Editorial note.} All delivered mathematical text is the author's own, copied from the pinned Stacks Project chapter file sites-cohomology.tex; no statement or proof step was rewritten, reordered or removed.\par\medskip
\noindent\textit{Editorial note.} The counted proof is the author's complete argument: the proof is written out in full: the choice of the K-injective resolution, the product complex representing the derived internal hom, the identification of the sheafified H\textasciicircum{}0 through the local comparison lemma, and the extension to all degrees giving the quasi-isomorphism.\par\medskip
\noindent\textit{Editorial note.} The label lemma-Rhom-strictly-perfect names the counted statement, and the proof environment immediately following it in the source is delivered verbatim as its proof.\par\medskip
\noindent\textit{Editorial note.} The counted statement is carried by the author's own lemma environment in the chapter source.\par\medskip
\noindent\textit{Editorial note.} The statement and proof are copied from the pinned chapter file; the two cited inputs are kept as external pointers with their canonical Stacks tags.\par\medskip
\noindent\textit{Editorial note.} The preamble is the chapter's own retained minimal preamble (sources/stacks-sites-cohomology/retained-preamble.tex), a verbatim excerpt of the native preamble: the Xy-pic setup and the native theorem declarations with the mathematical macros.\par\medskip
\noindent\textit{Allowed standard prerequisites (editorial).}\begin{itemize}
\item Strictly perfect complexes and the internal hom. Same chapter, the definition of a strictly perfect complex and the section constructing the internal hom complex with its differential. Supplies the hypotheses on E\textasciicircum{}bullet and the differential of the representing complex in the statement.
\item K-injective resolutions and the product complex representing R SheafHom. Same chapter, the construction representing R SheafHom by the product complex of internal homs into a K-injective resolution. Supplies the complex (H')\textasciicircum{}bullet used in the proof and its quasi-isomorphism to R SheafHom.
\item Local comparison of Hom sets for strictly perfect complexes. Same chapter, the lemma identifying the sheafification of the local Hom functor with the sheafification of the derived Hom for strictly perfect complexes. Supplies that the two complexes have the same sheafified cohomology, which is what makes them quasi-isomorphic.
\end{itemize}\medskip
\ifdefined\hypertarget\hypertarget{sourcepointer-1}{}\fi
\noindent\textit{Source reference.} Original reference 08C9: sources/stacks-sites-cohomology/pinned-original-source.tex, source lines 11881--11881. Prerequisite pointer invoked at proof line(s) 11950; it names a standard fact of the same chapter that the counted argument uses, and every step of the counted proof built on it is retained verbatim in the delivered body.\par\medskip
\ifdefined\hypertarget\hypertarget{sourcepointer-2}{}\fi
\noindent\textit{Source reference.} Original section 017G: sources/stacks-sites-cohomology/pinned-original-source.tex, source lines 8836--8836. Prerequisite pointer invoked at proof line(s) 11923; it names a standard fact of the same chapter that the counted argument uses, and every step of the counted proof built on it is retained verbatim in the delivered body.\par\medskip
\begin{lemma}
\label{lemma-Rhom-strictly-perfect}
Let $(\mathcal{C}, \mathcal{O})$ be a ringed site.
Let $\mathcal{E}^\bullet$, $\mathcal{F}^\bullet$ be complexes
of $\mathcal{O}$-modules with $\mathcal{E}^\bullet$ strictly perfect.
Then the internal hom $R\SheafHom(\mathcal{E}^\bullet, \mathcal{F}^\bullet)$
is represented by the complex $\mathcal{H}^\bullet$ with terms
$$
\mathcal{H}^n =
\bigoplus\nolimits_{n = p + q}
\SheafHom_\mathcal{O}(\mathcal{E}^{-q}, \mathcal{F}^p)
$$
and differential as described in Section \ref{section-internal-hom}.
\end{lemma}

\begin{proof}
Choose a quasi-isomorphism $\mathcal{F}^\bullet \to \mathcal{I}^\bullet$
into a K-injective complex. Let $(\mathcal{H}')^\bullet$ be the
complex with terms
$$
(\mathcal{H}')^n =
\prod\nolimits_{n = p + q}
\SheafHom_\mathcal{O}(\mathcal{L}^{-q}, \mathcal{I}^p)
$$
which represents $R\SheafHom(\mathcal{E}^\bullet, \mathcal{F}^\bullet)$
by the construction in Section \ref{section-internal-hom}. It suffices
to show that the map
$$
\mathcal{H}^\bullet \longrightarrow (\mathcal{H}')^\bullet
$$
is a quasi-isomorphism. Given an object $U$ of $\mathcal{C}$ we have
by inspection
$$
H^0(\mathcal{H}^\bullet(U)) =
\Hom_{K(\mathcal{O}_U)}(\mathcal{E}^\bullet|_U, \mathcal{I}^\bullet|_U)
\to
H^0((\mathcal{H}')^\bullet(U)) =
\Hom_{D(\mathcal{O}_U)}(\mathcal{E}^\bullet|_U, \mathcal{I}^\bullet|_U)
$$
By Lemma \ref{lemma-local-actual} the sheafification of
$U \mapsto H^0(\mathcal{H}^\bullet(U))$
is equal to the sheafification of
$U \mapsto H^0((\mathcal{H}')^\bullet(U))$. A similar argument can be
given for the other cohomology sheaves. Thus $\mathcal{H}^\bullet$
is quasi-isomorphic to $(\mathcal{H}')^\bullet$ which proves the lemma.
\end{proof}

\end{document}
